% Fiche 3 — Popper centralisé : la boucle et ses preuves
% Compilation : pdflatex fiche-popper-centralise.tex   (a besoin de style-fiche.tex)
\documentclass[10pt]{article}
% =====================================================================
%  Style commun aux fiches (bleu marine + bleu pastel)
%  Pour changer les couleurs : modifier uniquement ce fichier.
% =====================================================================
\usepackage[utf8]{inputenc}
\usepackage[T1]{fontenc}
\usepackage[scaled]{helvet}
\renewcommand{\familydefault}{\sfdefault}
\usepackage[a4paper,margin=8mm]{geometry}
\usepackage{xcolor}
\usepackage[most]{tcolorbox}
\usepackage{tikz}
\usepackage{pifont}
\usepackage{enumitem}
\usepackage{array}
\usepackage{etoolbox}
\usepackage{needspace}
\usepackage{amsmath,amssymb}
\usepackage{multirow}
\pagestyle{empty}
\setlength{\parindent}{0pt}

% ---------- Couleurs ----------
\definecolor{marine}{HTML}{3F72B8}   % bleu (plus clair que le marine)
\definecolor{pastel}{HTML}{BFD6F6}   % bleu pastel
\definecolor{fond}{HTML}{F4F7FB}
\definecolor{encre}{HTML}{1F1F1F}
\definecolor{vert}{HTML}{22A55B}
\definecolor{rouge}{HTML}{DC2626}
\pagecolor{fond}
\color{encre}

% ---------- Styles page de synthèse ----------
\newtcolorbox{bloc}[2][]{enhanced, colback=white, colframe=pastel,
  boxrule=0.9pt, arc=0pt, left=5pt, right=5pt, top=3pt, bottom=3pt,
  fontupper=\footnotesize\raggedright, #1,
  before upper={{\large\bfseries\textcolor{marine}{#2}}\par\vspace{2pt}}}
\newtcolorbox{astuce}{colback=pastel!75, colframe=pastel!75, arc=0pt,
  boxrule=0pt, left=4pt, right=4pt, top=2pt, bottom=2pt,
  fontupper=\footnotesize\raggedright}
\newcommand{\num}[1]{\textcolor{marine}{\textbf{#1.}}~}
\newcommand{\fl}{\tikz[baseline=-0.6ex]\draw[->,thick](0,0)--(0.4,0);}
\newcommand{\sous}[1]{\textcolor{marine}{\textbf{#1}}\par}
\setlist[itemize]{leftmargin=9pt, itemsep=0pt, topsep=1pt, parsep=0pt,
  label=\textbullet, before=\raggedright}
\newcolumntype{L}[1]{>{\raggedright\arraybackslash}p{#1}}

% Terme anglais (en gras) avec sa traduction française
\newcommand{\EN}[2]{\textbf{#1}\,{\scriptsize\textcolor{marine}{(#2)}}}

% ---------- Styles pages détaillées ----------
\newtcolorbox{dire}{enhanced, breakable, colback=pastel!25, colframe=marine,
  boxrule=0pt, leftrule=3pt, arc=0pt, left=8pt, right=6pt, top=4pt, bottom=4pt,
  fontupper=\small\itshape, before upper={{\upshape\bfseries\textcolor{marine}{FR}}\enspace}}
\newtcolorbox{direEN}{enhanced, breakable, colback=vert!6, colframe=vert,
  boxrule=0pt, leftrule=3pt, arc=0pt, left=8pt, right=6pt, top=4pt, bottom=4pt,
  fontupper=\small\itshape, before upper={{\upshape\bfseries\textcolor{vert}{EN}}\enspace}}
\newtcolorbox{pourquoi}{colback=fond, colframe=fond, arc=0pt, boxrule=0pt,
  left=6pt, right=6pt, top=3pt, bottom=3pt, fontupper=\small}
\newcommand{\titrebox}[2]{\needspace{10\baselineskip}\vspace{8pt}
  \begin{tcolorbox}[colback=marine, colframe=marine, arc=0pt, boxrule=0pt,
    top=4pt, bottom=4pt, left=8pt]
  \color{white}\Large\bfseries #1.~#2
  \end{tcolorbox}\vspace{2pt}}
\newcommand{\ce}[1]{\begin{pourquoi}\textbf{\textcolor{marine}{À retenir :}} #1\end{pourquoi}}
\newcommand{\st}[1]{\par\needspace{7\baselineskip}\vspace{5pt}{\normalsize\bfseries\textcolor{marine}{#1}}\par\vspace{2pt}}
\newcommand{\sst}[1]{\par\needspace{5\baselineskip}\vspace{3pt}{\small\bfseries\textcolor{encre}{#1}}\par\vspace{1pt}}
\newtcolorbox{moi}{colback=pastel!35, colframe=pastel!35, arc=0pt, boxrule=0pt,
  left=6pt, right=6pt, top=3pt, bottom=3pt, fontupper=\small}
\newtcolorbox{piege}{colback=rouge!6, colframe=rouge, boxrule=0pt, leftrule=3pt, arc=0pt,
  left=8pt, right=6pt, top=3pt, bottom=3pt, fontupper=\small}
% Boîte "énoncé" pour lemmes, propositions, théorèmes
\newtcolorbox{enonce}[1]{enhanced, breakable, colback=white, colframe=marine,
  boxrule=0.8pt, arc=0pt, left=8pt, right=8pt, top=4pt, bottom=4pt,
  fontupper=\small, title={#1}, coltitle=white, colbacktitle=marine,
  fonttitle=\bfseries\small}
% Passage à la mise en page détaillée
\newcommand{\detailpages}{%
  \newpage
  \newgeometry{a4paper,margin=15mm}
  \setlist[itemize]{leftmargin=12pt, itemsep=1pt, topsep=2pt, parsep=0pt,
    label=\textcolor{marine}{\textbullet}}
  \small}


\begin{document}

% ================= EN-TÊTE =================
\begin{minipage}[c]{0.64\textwidth}
  {\fontsize{26}{28}\selectfont\bfseries\textcolor{marine}{PROOF}~CHEAT SHEET 3}\\[1pt]
  {\small Popper (centralized) --- generate, test, constrain : pourquoi ça marche}
\end{minipage}\hfill
\begin{minipage}[c]{0.34\textwidth}\raggedright
  \normalsize \textbf{Sound pruning} $\to$ \textbf{termination}\\ $\to$ \textbf{finds a solution if one exists}
\end{minipage}

\vspace{5pt}

% ================= LIGNE 1 =================
\begin{minipage}[t]{0.32\textwidth}
\begin{bloc}{\num{1}The loop (la boucle)}
\begin{enumerate}[leftmargin=12pt, itemsep=1pt, topsep=1pt]
  \item \textbf{Generate} (ASP, clingo) : un programme $H_t$ qui respecte toutes les contraintes $\mathcal{C}_{t-1}$.
  \item \textbf{Test} (Prolog) : TP, FN, FP, TN $\to$ issue $(\varepsilon^+,\varepsilon^-)$.
  \item \textbf{Constrain} : de $(H_t,\text{issue})$ on tire des contraintes qui \emph{interdisent} des familles de programmes.
\end{enumerate}
On s'arrête quand l'issue est (ALL, NONE) : $H_t$ est une \textbf{solution} (complète et cohérente).
\begin{astuce}
Le solveur est chargé \textbf{une fois} (multi-shot) : on lui ajoute des contraintes à chaque tour.
\end{astuce}
\end{bloc}
\end{minipage}\hfill
\begin{minipage}[t]{0.32\textwidth}
\begin{bloc}{\num{2}Outcome $\to$ constraints}
{\scriptsize
\renewcommand{\arraystretch}{1.15}
\begin{tabular}{@{}l l L{2.6cm}@{}}
$\varepsilon^+$ & $\varepsilon^-$ & Prunes (élague)\\\hline
NONE & NONE & Spec + Redundancy\\
NONE & SOME & Spec + Red. + Gen\\
SOME & NONE & Spec\\
SOME & SOME & Spec + Gen\\
ALL & SOME & Gen\\
ALL & NONE & \textbf{Solution}\\
\end{tabular}}
\vspace{2pt}
\begin{itemize}
  \item \textbf{Spec} si $\varepsilon^+\neq\mathit{ALL}$ (positifs manqués)
  \item \textbf{Gen} si $\varepsilon^-=\mathit{SOME}$ (négatifs couverts)
\end{itemize}
\end{bloc}
\end{minipage}\hfill
\begin{minipage}[t]{0.32\textwidth}
\begin{bloc}{\num{3}Three properties}
\begin{itemize}
  \item \textbf{Sound pruning} : une contrainte n'interdit jamais une solution.
  \item \textbf{Termination} : au plus $|\mathcal{H}|$ tours (chaque tour interdit au moins $H_t$).
  \item \textbf{Completeness of the search} : s'il existe une solution dans $\mathcal{H}$, Popper en renvoie une.
\end{itemize}
\begin{astuce}
Tout repose sur la monotonie : plus de clauses $\Rightarrow$ au moins autant de couverture.
\end{astuce}
\end{bloc}
\end{minipage}

\vspace{4pt}

% ================= LIGNE 2 =================
\begin{minipage}[t]{0.5\textwidth}
\begin{bloc}{\num{4}Lemma : sound pruning}
\textbf{Définitions.} $H'$ est une \textbf{generalisation} de $H$ si $H'\models H$ ; une \textbf{specialisation} de $H$ si $H\models H'$.\par\vspace{2pt}
\textbf{(a)} Si $H$ couvre un négatif $e^-$, toute généralisation $H'$ le couvre aussi $\Rightarrow$ $H'$ est \emph{incohérente}, pas une solution.\par\vspace{2pt}
\textbf{(b)} Si $H$ rate un positif $e^+$, toute spécialisation $H'$ le rate aussi $\Rightarrow$ $H'$ est \emph{incomplète}, pas une solution.\par\vspace{2pt}
\textbf{Preuve} : $H'\models H$ et $\mathcal{B}\cup H\models e$ donnent $\mathcal{B}\cup H'\models e$ (a). $H\models H'$ et $\mathcal{B}\cup H'\models e$ donnent $\mathcal{B}\cup H\models e$ (b, par contraposée).
\end{bloc}
\end{minipage}\hfill
\begin{minipage}[t]{0.48\textwidth}
\begin{bloc}{\num{5}Theorems}
\textbf{Termination.} Chaque tour, $H_t$ est interdit (il est sa propre généralisation et spécialisation). Donc $\le|\mathcal{H}|$ tours.\par\vspace{3pt}
\textbf{Completeness of the search.} Soit $S$ une solution de $\mathcal{H}$. Par le lemme, $S$ n'est jamais interdite. Donc le générateur ne renvoie jamais $\bot$. Comme la boucle termine, elle s'arrête sur le test (ALL, NONE) : elle renvoie une solution.\par\vspace{3pt}
\textbf{Size-minimality} (si le générateur procède par taille croissante) : les tailles inférieures sont épuisées sans solution, donc la solution renvoyée est de taille minimale.
\end{bloc}
\end{minipage}

\vspace{4pt}

% ================= LIGNE 3 =================
\begin{minipage}[t]{0.52\textwidth}
\begin{bloc}{\num{6}On Trains (centralisé)}
{\scriptsize
\renewcommand{\arraystretch}{1.1}
\begin{tabular}{@{}l l@{}}
Initial space & 1433 programs (sizes 3--6)\\
By size & 6 / 29 / 185 / 1213\\
Rounds & 48 (47 failures + 1 solution)\\
Solution & size 6, found at round 48\\
\end{tabular}}\par\vspace{2pt}
\footnotesize Seuls 48 programmes sur 1433 sont testés : les contraintes élaguent tout le reste sans le tester.
\end{bloc}
\end{minipage}\hfill
\begin{minipage}[t]{0.46\textwidth}
\begin{bloc}{\num{7}Pitfalls (les pièges)}
\begin{itemize}
  \item \textbf{Les contraintes ne disent pas ``va par là''} : elles disent ``pas par là''. Les spécialisations d'un programme trop général restent possibles.
  \item \textbf{Redundancy} : vient d'une analyse séparée (Morel et al.), pas du lemme ci-contre.
  \item \textbf{Implémentation} : Popper n'élague qu'une partie sûre des familles (par exemple $\theta$ injectif, spécialisations à une clause).
  \item \textbf{Bruit} : la preuve suppose des exemples sans bruit.
\end{itemize}
\end{bloc}
\end{minipage}

\vspace{5pt}

\begin{tcolorbox}[colback=marine, colframe=marine, arc=0pt, boxrule=0pt,
  halign=center, top=3pt, bottom=3pt]
\color{white}\small Popper never prunes a solution, and it prunes at least one program per round: it must find a solution if there is one.
\end{tcolorbox}

% =====================================================================
%                    PAGES DÉTAILLÉES
% =====================================================================
\detailpages

{\fontsize{24}{26}\selectfont\bfseries\textcolor{marine}{POPPER}~CENTRALISÉ}\\[2pt]
{\normalsize La boucle \EN{generate-test-constrain}{générer-tester-contraindre} et les preuves qui expliquent pourquoi elle est correcte.
Les termes de théorie sont donnés en anglais (en \textbf{gras}).}

% ---------------------------------------------------------------- 1
\titrebox{1}{The ILP problem (le problème)}
Un problème ILP est un quadruplet $(\mathcal{B},E^+,E^-,\mathcal{H})$ :
\begin{itemize}
  \item $\mathcal{B}$ : le \EN{background knowledge}{connaissances de fond}, un ensemble fini de faits et de règles (sur Trains : \texttt{has\_car}, \texttt{has\_load}, \texttt{short}, \texttt{long}\ldots).
  \item $E^+$, $E^-$ : \EN{positive and negative examples}{exemples positifs et négatifs}, ici 5 et 5 atomes \texttt{f(t\ldots)}.
  \item $\mathcal{H}$ : l'\EN{hypothesis space}{espace d'hypothèses}, fini, défini par le \EN{language bias}{biais de langage}
        (\texttt{bias.pl} : \texttt{max\_clauses(3)}, \texttt{max\_vars(5)}, \texttt{max\_body(6)}, types, directions).
\end{itemize}
\textbf{Couverture} (\EN{coverage}{couverture}) : $H$ couvre $e$ si $\mathcal{B}\cup H\models e$.\par\vspace{2pt}
\textbf{Solution} : une hypothèse $H\in\mathcal{H}$ qui est \EN{complete}{complète} (couvre tous les positifs) et \EN{consistent}{cohérente} (ne couvre aucun négatif).
C'est exactement l'issue (ALL, NONE).

% ---------------------------------------------------------------- 2
\titrebox{2}{The three steps (les trois étapes)}

\st{Generate : le générateur ASP}
Popper charge dans \texttt{clingo} trois morceaux : \texttt{alan.pl} (ce qu'est un programme bien formé, identique pour tous les problèmes),
\texttt{bias.pl} (les faits du problème) et les contraintes apprises. Un \EN{model}{modèle} ASP correspond à un programme valide.
On avance par taille croissante avec \texttt{size\_in\_literals(n)}. Le solveur est en \EN{multi-shot}{multi-coup} : il est chargé une fois, et les contraintes y sont ajoutées à chaque tour.\par\vspace{2pt}
Sur Trains : espace initial de 1433 programmes (6 / 29 / 185 / 1213 aux tailles 3 à 6).

\st{Test : Prolog}
Le programme est chargé avec le BK, et pour chaque exemple Prolog essaie de prouver le but (résolution SLD, fiche~1), avec un \EN{timeout}{délai maximal}.
On obtient TP, FN, FP, TN puis l'issue $(\varepsilon^+,\varepsilon^-)$.

\st{Constrain : les contraintes}
De $(H,\text{issue})$ on dérive des contraintes ASP qui \EN{prune}{élaguent} l'espace. Elles sont instanciées (\EN{grounding}{instanciation}) puis ajoutées au solveur.

% ---------------------------------------------------------------- 3
\titrebox{3}{Outcome and constraints (issue et contraintes)}
\renewcommand{\arraystretch}{1.25}
\begin{tabular}{@{}l l L{5.2cm} L{5.2cm}@{}}
$\varepsilon^+$ & $\varepsilon^-$ & \textbf{Constraints derived} & \textbf{Why}\\\hline
NONE & NONE & Specialisation + Redundancy & rien n'est couvert : les spécialisations ne couvrent pas plus\\
NONE & SOME & Specialisation + Redundancy + Generalisation & rien de positif, mais des négatifs couverts\\
SOME & NONE & Specialisation & incomplet, cohérent\\
SOME & SOME & Specialisation + Generalisation & incomplet et incohérent\\
ALL & SOME & Generalisation & complet, incohérent\\
ALL & NONE & \textbf{Solution found} & on s'arrête\\
\end{tabular}
\ce{\textbf{Spec} dès que $\varepsilon^+\neq\mathit{ALL}$ ; \textbf{Gen} dès que $\varepsilon^-=\mathit{SOME}$. Les contraintes disent ``pas par là'', jamais ``va par là''.}

\st{Un exemple réel : le programme \#1 \texttt{f(A) :- has\_car(A,B), short(B).}}
Il couvre une partie des positifs et des négatifs : issue (SOME, SOME) (comptes détaillés dans la fiche~2). Popper ajoute alors des règles ASP instanciées, dont une définition et deux interdictions (forme simplifiée) :
\begin{moi}\footnotesize\ttfamily
included\_clause(fAhas\_carABshortB,Cl) :- head\_literal(Cl,f,1,(A,)),\\
\hspace*{2em} body\_literal(Cl,short,1,(B,)), body\_literal(Cl,has\_car,2,(A,B)), \ldots\\[2pt]
\% generalisation\\
:- included\_clause(fAhas\_carABshortB,C0), body\_size(C0,2), \ldots\\
\% specialisation\\
:- included\_clause(fAhas\_carABshortB,C0), \ldots, not clause(1)
\end{moi}
La règle \texttt{included\_clause} est un \textbf{détecteur} (une définition, sans \texttt{:-} en tête). Les lignes qui commencent par \texttt{:-} sont les \textbf{interdictions}.

% ---------------------------------------------------------------- 4
\titrebox{4}{Generalisation and specialisation (via entailment)}
\begin{enonce}{Definition (Generalisation, Specialisation)}
Soient $H,H'$ deux hypothèses. $H'$ est une \EN{generalisation}{généralisation} de $H$ si $H'\models H$.
$H'$ est une \EN{specialisation}{spécialisation} de $H$ si $H\models H'$.
\end{enonce}
Intuition : ajouter une clause à un programme (le \emph{généraliser}) ne retire aucune conséquence ; ajouter un littéral dans une clause (la \emph{spécialiser})
n'en ajoute pas.\par\vspace{2pt}
\textbf{Mise en garde.} Vérifie la formulation exacte dans l'article de Cropper et Morel (Popper, 2021) : certaines présentations passent par la
\EN{$\theta$-subsumption}{$\theta$-subsomption} ($c_1\theta\subseteq c_2$). L'implémentation de Popper utilise une version plus stricte (\texttt{AllDifferent} : $\theta$ injectif sur les variables).
Les deux sont \emph{sûres} : elles n'élaguent qu'un sous-ensemble de ce que la définition autorise.

% ---------------------------------------------------------------- 5
\titrebox{5}{Lemma 1 (monotonicity of coverage)}
\begin{enonce}{Lemma 1}
Soit $e$ un exemple. (a) Si $H'$ est une généralisation de $H$ et $\mathcal{B}\cup H\models e$, alors $\mathcal{B}\cup H'\models e$.
(b) Si $H'$ est une spécialisation de $H$ et $\mathcal{B}\cup H'\models e$, alors $\mathcal{B}\cup H\models e$.
\end{enonce}
\textbf{Proof.} (a) $H'\models H$, donc tout modèle de $\mathcal{B}\cup H'$ est un modèle de $\mathcal{B}\cup H$, qui satisfait $e$.
(b) $H\models H'$, donc tout modèle de $\mathcal{B}\cup H$ est un modèle de $\mathcal{B}\cup H'$, qui satisfait $e$. \hfill $\square$

% ---------------------------------------------------------------- 6
\titrebox{6}{Lemma 2 (sound pruning)}
\begin{enonce}{Lemma 2 (Sound Pruning)}
Soit $H$ un programme testé. (a) Si $\varepsilon^-(H)=\mathit{SOME}$, toute généralisation de $H$ est incohérente. (b) Si $\varepsilon^+(H)\neq\mathit{ALL}$, toute spécialisation de $H$ est incomplète.
Donc aucune solution n'est jamais élaguée par une contrainte de généralisation ou de spécialisation.
\end{enonce}
\textbf{Proof.} (a) Il existe un négatif $e^-$ avec $\mathcal{B}\cup H\models e^-$ ; par le Lemme 1(a), toute généralisation $H'$ de $H$ vérifie $\mathcal{B}\cup H'\models e^-$, donc couvre un négatif.
(b) Il existe un positif $e^+$ avec $\mathcal{B}\cup H\not\models e^+$ ; si une spécialisation $H'$ de $H$ vérifiait $\mathcal{B}\cup H'\models e^+$, le Lemme 1(b) donnerait $\mathcal{B}\cup H\models e^+$, contradiction. Donc $H'$ rate $e^+$ : elle est incomplète.
Une solution est complète et cohérente : elle ne peut être ni dans (a) ni dans (b). \hfill $\square$

\st{Pourquoi la table des contraintes est cohérente avec ce lemme}
\begin{itemize}
  \item $\varepsilon^+\in\{\mathit{NONE},\mathit{SOME}\}$ $\Rightarrow$ $\varepsilon^+\neq\mathit{ALL}$ $\Rightarrow$ contrainte de \textbf{spécialisation}.
  \item $\varepsilon^-=\mathit{SOME}$ $\Rightarrow$ contrainte de \textbf{généralisation}.
  \item Les contraintes de \textbf{redundancy} (cas $\varepsilon^+=\mathit{NONE}$) viennent d'une autre analyse (Morel et al.) : \emph{elles ne sont pas prouvées ici}.
\end{itemize}

% ---------------------------------------------------------------- 7
\titrebox{7}{Theorems (termination, completeness of the search)}
\begin{enonce}{Theorem 1 (Termination)}
La boucle de Popper s'arrête après au plus $|\mathcal{H}|$ tours.
\end{enonce}
\textbf{Proof.} À chaque tour non terminal, l'issue n'est pas (ALL, NONE), donc $\varepsilon^+\neq\mathit{ALL}$ ou $\varepsilon^-=\mathit{SOME}$. Dans le premier cas la contrainte de spécialisation interdit $H_t$ (qui est sa propre spécialisation, car $H_t\models H_t$) ; dans le second, celle de généralisation l'interdit. Le générateur ne propose que des programmes non interdits : $H_t$ ne reviendra pas. Comme $\mathcal{H}$ est fini, il y a au plus $|\mathcal{H}|$ tours. \hfill $\square$

\begin{enonce}{Theorem 2 (Completeness of the search)}
(i) Si Popper renvoie $H$, alors $H$ est une solution. (ii) S'il existe une solution dans $\mathcal{H}$, Popper en renvoie une.
\end{enonce}
\textbf{Proof.} (i) Popper ne renvoie que quand l'issue est (ALL, NONE), c'est-à-dire quand $H$ est complète et cohérente.
(ii) Soit $S$ une solution. Par le Lemme 2, $S$ n'est jamais élaguée ; elle reste donc disponible pour le générateur, qui ne renvoie jamais $\bot$. Par le Théorème 1 la boucle termine ; la seule façon de terminer sans $\bot$ est de renvoyer sur l'issue (ALL, NONE), c'est-à-dire de renvoyer une solution (pas forcément $S$). \hfill $\square$

\begin{moi}
\textbf{Remarque (taille minimale).} Si le générateur propose les programmes par taille non décroissante, tous les programmes de taille inférieure à celle de la solution renvoyée ont été testés ou élagués, et, par le Lemme 2, aucun d'eux n'était une solution : la solution renvoyée est de \textbf{taille minimale}. Sur Trains, les tailles 3, 4 et 5 sont épuisées avant la solution de taille 6.
\end{moi}

% ---------------------------------------------------------------- 8
\titrebox{8}{Lien avec FedPopper}
\begin{itemize}
  \item Popper centralisé = Algorithme~1 de FedPopper avec \textbf{un seul client} qui détient le dataset poolé.
  \item Les preuves de cette fiche (lemme et théorèmes) \textbf{ne changent pas} en fédéré : elles portent sur $(H_t,\text{issue},\mathcal{C}_t)$.
  \item La fiche~2 montre que FedPopper calcule \textbf{la même issue} à chaque tour, donc hérite de ces propriétés.
\end{itemize}

% ---------------------------------------------------------------- 9
\titrebox{9}{Trains : la trace complète}
\begin{tabular}{@{}l r@{}}
\textbf{Size of programs} & \textbf{Programs in initial space}\\\hline
3 & 6\\
4 & 29\\
5 & 185\\
6 & 1213\\
\textbf{Total} & \textbf{1433}\\
\end{tabular}
\ce{48 programmes testés (47 échecs puis la solution) : tout le reste de l'espace est élagué par les contraintes. La solution est de taille 6. Le détail par taille des programmes testés n'est pas reproduit ici : à relever dans le log de Popper si tu veux le citer.}
\textbf{Solution :} \texttt{f(A) :- has\_car(A,B), has\_car(A,C), long(C), three\_wheels(B), two\_wheels(C).}
(5 positifs couverts, aucun négatif.)

% ---------------------------------------------------------------- 10
\titrebox{10}{Reviewer questions (questions de relecteur)}

\st{``Popper peut-il éliminer une solution par erreur ?''}
\begin{dire}
Non : le lemme d'élagage sûr montre que les contraintes de généralisation et de spécialisation n'éliminent que des programmes incohérents ou incomplets.
\end{dire}
\begin{direEN}
No: the sound-pruning lemma shows that generalisation and specialisation constraints only remove programs that are inconsistent or incomplete.
\end{direEN}

\st{``Pourquoi la boucle termine-t-elle ?''}
\begin{dire}
L'espace d'hypothèses est fini et chaque tour interdit au moins l'hypothèse testée.
\end{dire}
\begin{direEN}
The hypothesis space is finite and each round forbids at least the tested hypothesis.
\end{direEN}

\st{``Et avec du bruit dans les données ?''}
\begin{dire}
Les preuves supposent des exemples sans bruit : un positif mal étiqueté rend l'élagage par spécialisation faux. Le cas bruité demande un autre critère (par exemple MDL).
\end{dire}
\begin{direEN}
The proofs assume noise-free examples: a mislabeled positive makes specialisation pruning unsound. The noisy case needs another criterion (for instance MDL).
\end{direEN}

\vspace{8pt}
\begin{tcolorbox}[colback=marine, colframe=marine, arc=0pt, boxrule=0pt,
  halign=center, top=4pt, bottom=4pt]
\color{white}\small Never prune a solution + prune at least one program per round = terminate with a solution if one exists.
\end{tcolorbox}

\end{document}
